\documentclass[../../infufrgs/ppgc-prop-tese.tex]{subfiles}
\begin{document}


%It is sad to announce that besides a steering method and an in depth analysis of late generalization this work does not entail a new method, or criteria for making models work better, faster or safer. Nonetherless we provide an analysis if interesting, sometimes surprising, behavior of such models and insights. It is a collection of insights about the inner-workings of attention-based models that contribute to a better understaning and appreciation of their emergent phenomena.

This thesis builds on results obtained during the PhD program, including an accepted conference paper and consolidated experimental pipelines. The main completed outcomes are:
%
\begin{enumerate*}[label=(\roman*)]
  \item a set of controlled experiments and analyses on late generalization (grokking-like transitions), consolidated in the paper “Inducing Grokking with Distribution Shifts”;
  \item an implemented prototype of a logic-inspired steering/constraint framework, with preliminary experiments that inform the final methodology and ablation plan; and
  \item an initial set of observations on internal activation/attention patterns under interventions and input perturbations, motivating the diagnostic component of the thesis.
\end{enumerate*}
%
The remaining work consists of consolidating the final steering method, running definitive experiments with baselines and ablations, and integrating all results into a coherent thesis narrative with reproducibility artifacts.
\end{document}